\documentclass[12pt]{report}
\usepackage[italian]{babel}
\usepackage[letterpaper,top=2cm,bottom=2cm,left=3cm,right=3cm,marginparwidth=1.75cm]{geometry}
\usepackage{amsmath}
\usepackage{siunitx}
\usepackage{graphicx}
\usepackage[colorlinks=true, allcolors=blue]{hyperref}
\usepackage{soul}
\usepackage{xcolor}
\sethlcolor{yellow}
\usepackage{titlesec}
\usepackage{float}
\usepackage[
  backend=biber,
  style=numeric, 
  sorting=none,         
  maxnames=3,
  doi=true,
  url=true
  ]{biblatex}  
\addbibresource{references.bib}
\nocite{*}
\begin{document}
\begin{titlepage}
    \centering
    
    \begin{center}
        \includegraphics[width=8cm]{LogoReport.png}
    \end{center}

    \vspace{2cm}

    {\Huge\bfseries Report\par}

    \vspace{0.5cm}

    {\Large Dinamica e controllo RITM-200\par}

    \vfill

    \rule{0.8\textwidth}{0.6pt}

    \vspace{0.8cm}

    {\Large\textbf{Autori}}

    \vspace{0.5cm}

    \begin{tabular}{c}
        Ettore Carlotti\\
        Andrea Ferroni\\
        Gabriele Melchionda\\
        Vsevolod Ozmitel Volpato\\
        Giuseppe Praino
    \end{tabular}

    \vspace*{\fill}

    \large
    Laboratorio di ingegneria nucleare\\
    Anno Accademico 2025 - 2026\\
    \today

\end{titlepage}

\begin{abstract}
Il presente lavoro analizza il comportamento dinamico del reattore nucleare 
di tipo SMR RITM-200 in seguito a perturbazioni operative, mediante la 
realizzazione di un modello numerico accoppiato implementato in Python.\\\\
Il modello integra tre sottosistemi fisici: la cinetica neutronica puntiforme 
a sei gruppi di neutroni ritardati, la termoidraulica del nocciolo con tre nodi radiali per il combustibile, e un modello monodimensionale 
a volumi finiti per lo scambiatore di calore a tubi concentrici in controcorrente. 
La dinamica della pressione del circuito secondario è descritta tramite un modello 
fenomenologico del primo ordine, ricavato dal bilancio di energia interna del 
fluido secondario.\\\\
L'integrazione temporale del sistema accoppiato è realizzata con uno schema 
implicito di Eulero per lo scambiatore di calore, incondizionatamente stabile, 
e con il metodo \texttt{LSODA} per le equazioni differenziali del nocciolo, 
adatto a sistemi stiff. I due modelli sono accoppiati attraverso le temperature 
di \textit{hot leg} e \textit{cold leg} del circuito primario, aggiornate ad 
ogni passo temporale.\\\\
Sono analizzati due scenari transitori: un aumento della portata 
massica del circuito secondario, che illustra la capacità di \textit{load 
following} caratteristica dei reattori PWR, e un inserimento delle barre di 
controllo nel nocciolo, che dimostra il controllo della potenza a monte. 
I risultati mostrano in entrambi i casi la corretta propagazione degli effetti 
tra i sottosistemi e la stabilità del nuovo stato di equilibrio raggiunto, 
confermando la validità del modello sviluppato per l'analisi di transitori 
operativi in reattori SMR.
\end{abstract}

\tableofcontents

\chapter{Introduzione}

\section{Descrizione del reattore}
Il RITM-200 è un reattore di tipo SMR di generazione III+ progettato per produrre fino a 55MWe secondo tecnologia PWR. Ha un'aspettativa di vita di 60 anni, il fuel utilizzato è uranio arricchito al 14,06\%, il cui ciclo termodinamico dipende dall'applicazione del reattore (propulsione navale artica, produzione elettrica su centrali galleggianti o future centrali modulari terrestri), richiedendo refueling ogni 6 anni per applicazioni terrestri o 10 anni per la propulsione navale. Essendo inizialmente progettato per navi rompighiaccio presenta una struttura compatta e leggera, tutti i componenti principali sono infatti incorporati nel vessel: il nocciolo si trova nella parte centrale inferiore, i 4 scambiatori di calore sono disposti radialmente attorno al nocciolo e la pressione viene regolata nella parte superiore del vessel senza necessità di avere un pressurizzatore. Questa configurazione differisce da quella tipica delle grandi centrali PWR, la riduzione di tubazioni presenta un vantaggio non solo economico ma anche dal punto di vista della sicurezza.\\
\begin{center}
\begin{table}[htbp]
    \centering
    \caption{Valori operativi}
    \label{1.1}

\begin {tabular}{|m{8cm}|m{3cm}|}
    \hline
    Pressione  operativa primario & $15.7\, \text{MPa}$\\
    \hline
    Temperatura di ingresso refrigerante & $277\, \text{°C}$\\
    \hline
    Temperatura in uscita refrigerante & 313 °C\\
    \hline
    Portata massica del refrigerante & 3250 t/h\\
    \hline
    Potenza termica nominale & 175 MW\\
    \hline
    Potenza elettrica & $ 2 \times 55\,\text {MW}$\\
    \hline
    Pressione secondario & $3.4\, \text{MPa}$\\
    \hline

\end{tabular}
\end{table}
\end {center}

\section{Metodi numerici}

La risoluzione del problema è stata implementata in Python, facilitando la risoluzione di sistemi di equazioni differenziali ordinarie di natura stiff (ossia in cui due variabili evolvono su scale temporali molto diverse tra loro), sistemi non lineari per le condizioni stazionarie e discretizzazione spaziale delle equazioni di trasporto nello scambiatore di calore. Il codice impiega le librerie \texttt{numPy} e \texttt{sciPy}, insieme alla libreria \texttt{pyXSteam} per il calcolo delle proprietà dell'acqua. 
%\texttt{numpy} viene utilizzato per le operazioni vettoriali e matriciali; \texttt{scipy} fornisce le funzioni utili per risolvere i problemi numerici: \texttt{solve\_ivp} per l'integrazione temporale dei sistemi di ODE, \texttt{fsolve} per la soluzione dei sistemi algebrici non lineari in regime stazionario, e \texttt{root\_scalar} per la ricerca dello zero nel ciclo di accoppiamento core-SG. Le proprietà termodinamiche dell'acqua e del vapore sono ottenute dalla libreria \texttt{pyXSteam}, che implementa le tavole IF-97. Infine, \texttt{matplotlib.pyplot} è stata utile per il tracciamento dei grafici che forniscono una rappresentazione visiva dei risultati.

%L'integrazione temporale adotta il metodo LSODA, che permette di avere soluzioni veloci anche nel caso dei sistemi "stiff", ossia in cui due variabili evolvono su scale temporali molto diverse tra loro. Infatti, i precursori di neutroni ritardati evolvono su scale di $10^{-2}$--$10^{0}$ s, mentre la risposta termica del combustibile e del coolant si sviluppa su scale di $10^{0}$--$10^{1}$ s. Il generatore di vapore è integrato con uno schema implicito a volumi finiti, incondizionatamente stabile, che non impone vincoli sul passo temporale di tipo CFL e consente quindi di usare $\Delta t = 0.01$ s per l'intero transitorio.

\subsection{Integrazione delle ODE del nocciolo}
Le equazioni del nocciolo combinano la cinetica neutronica a sei famiglie di neutroni ritardati con un modello termoidraulico a tre nodi radiali per il combustibile, il rivestimento e il refrigerante, per un totale di $6 + 5 = 11$ equazioni differenziali accoppiate. Il sistema presenta un'elevata rigidità (\textit{stiffness}) dovuta alla coesistenza di costanti di tempo molto diverse:
il tempo di vita dei neutroni prompt $\Lambda = 10^{-4}$~s contro le costanti termiche del combustibile e del refrigerante, dell'ordine di decine di secondi. Per far fronte a questa rigidità si impiega la funzione \texttt{solve\_ivp} di \texttt{scipy.integrate} con il metodo \texttt{LSODA} (Livermore Solver for ODEs with Automatic stiffness detection and method switching), che commuta
dinamicamente tra un integratore Adams-Moulton (adatto a problemi non rigidi) e un metodo BDF (Backward Differentiation Formula, implicito, del primo-quinto ordine) in funzione del numero di condizionamento locale della Jacobiana. Le tolleranze di integrazione sono impostate a $\texttt{rtol} = 10^{-6}$ e $\texttt{atol} = 10^{-8}$, garantendo precisione numerica adeguata senza oneri computazionali eccessivi.

\subsection{Soluzione dello stato stazionario accoppiato}

La ricerca dello stato stazionario accoppiato nocciolo--generatore di vapore avviene mediante due livelli di iterazione. Al livello interno, il sistema non lineare di cinque equazioni algebriche per le temperature del combustibile ($T_1$, $T_2$, $T_3$), del rivestimento ($T_{cl}$) e del refrigerante ($T_p$)
viene risolto con la funzione \texttt{fsolve} di \texttt{scipy.optimize}, che implementa un metodo ibrido di tipo Powell basato su differenze finite per
l'approssimazione della Jacobiana. Al livello esterno, la temperatura della cold leg che bilancia contemporaneamente il nocciolo e lo scambiatore viene ricercata con \texttt{root\_scalar} (sempre da \texttt{scipy.optimize}) applicando il metodo di Brent-Dekker (\texttt{method='brentq'}), che garantisce la convergenza in un numero finito di passi su un intervallo iniziale definito tra $T_{\text{s,in}}$ e $T_\text{hot}$.

\subsection{Schema implicito per il generatore di vapore}

Il generatore di vapore è modellato come uno scambiatore a tubi coassiali in controcorrente, discretizzato assialmente in $N = 30$ volumi di controllo. L'integrazione temporale delle equazioni di bilancio energetico per il primario e il secondario è eseguita con uno schema \emph{implicito di Eulero} su ogni cella finita, che, essendo incondizionatamente stabile, non presenta vincoli di tipo \textit{CFL}. La struttura \textit{upwind} del trasporto in controcorrente consente di fattorizzare il problema in due \textit{sweep} sequenziali privi di inversione matriciale: ad ogni passo la temperatura aggiornata della cella a valle è già disponibile, riducendo il sistema a una semplice sostituzione in avanti cella per cella.

\subsection{Procedura di warmup}
Inizialmente viene eseguita una fase di \textit{warmup} implicito di 1000 passi con $\Delta t_{wu} = 1$~s, volta a portare i profili discreti dello scambiatore allo stato stazionario coerente con lo schema numerico adottato. Ad ogni passo del warmup si esegue un ciclo con predizione (guess) e successiva correzione sull'accoppiamento \textit{SG}--core: data una stima della temperatura della \textit{cold leg} dal profilo corrente del primario, il correttore la aggiorna usando il profilo \textit{SG} appena integrato, e il core viene integrato con \texttt{solve\_ivp}/\texttt{LSODA} sull'intervallo $[0,\,\Delta t_{wu}]$. Poiché lo schema implicito non è soggetto a vincoli di stabilità \textit{CFL}, il passo $\Delta t_{wu}$ può essere scelto pari a $\sim 10\times$ il tempo di transito del refrigerante primario nel generatore di vapore, riducendo il numero di passi necessari alla convergenza di oltre due ordini di grandezza rispetto a uno schema esplicito.

\chapter{Neutronica}

In questo capitolo analizzeremo la cinetica neutronica del reattore. L'analisi neutronica, in questo contesto, permette di simulare la risposta temporale della densità neutronica, da cui dipende direttamente l'evoluzione della potenza termica generata nel nocciolo. Poiché il reattore in analisi, essendo di tipo SMR, è progettato per operare con elevati standard di stabilità e sicurezza passiva, è necessario studiare come il sistema reagisce dinamicamente a improvvise perturbazioni esterne o operative. Il modello si concentra sul calcolo della reattività totale del sistema ($\rho$) e sulla risoluzione delle equazioni differenziali della cinetica puntiforme. La reattività è stata espressa valutando i coefficienti di retroazione o feedback legati alle variazioni di temperatura del combustibile ($\alpha_f$), del moderatore ($\alpha_m$) e all'inserzione delle barre di controllo ($\alpha_h$). La modellazione spaziale tridimensionale del flusso neutronico è stata tuttavia trascurata, poiché l'approccio a zero dimensioni (0D), accoppiato all'uso di sei gruppi di precursori di neutroni ritardati, cattura con sufficiente accuratezza l'andamento transitorio globale necessario per questa analisi.

\section{ Fenomenologia fisica e retroazione}

 Il comportamento dinamico di un reattore nucleare, in particolare nel caso di sistemi compatti come il RITM-200, è governato da una stretta interazione non lineare tra la popolazione neutronica e le grandezze termoidrauliche che caratterizzano il nocciolo. Il principio cardine che regola la stabilità globale del sistema è il concetto di feedback della reattività. La reattività rappresenta la misura dello scostamento del reattore dalla condizione di criticità stazionaria: variazioni termiche, geometriche o meccaniche all'interno del core alterano istantaneamente l'equilibrio tra i neutroni prodotti per fissione e quelli persi per cattura o leakage, modificando di conseguenza la densità neutronica e la potenza termica generata.

\subsection{Meccanismi di feedback}

I meccanismi che governano questo equilibrio si dividono in contributi intrinseci termofisici del reattore, ed esogeni, cioè di controllo:

\subsubsection{\textbf{Meccanismi di feedback intrinseci}}

\subsubsection{Effetto Doppler del combustibile} Costituisce il principale meccanismo di sicurezza passiva: l’innalzamento di temperatura del combustibile aumenta l’agitazione termica dei nuclei del combustibile, allargandone le finestre energetiche di assorbimento (aumentando lo spettro di risonanza), a causa del moto relativo rispetto ai neutroni incidenti. Questo allargamento incrementa la cattura parassita di neutroni nel nocciolo, generando un’immediata introduzione di reattività negativa che autosmorza e stabilizza il reattore in modo istantaneo in caso di picchi di potenza.


\subsubsection{Feedback del Moderatore} L'aumento di temperatura dell'acqua ne provoca l'espansione e la conseguente riduzione di densità. Una minore densità degrada la capacità moderante del fluido: i neutroni mantengono energie più elevate, aumentando la probabilità di leakage o di cattura prima della termalizzazione. Ciò introduce reattività negativa. Viceversa, un raffreddamento improvviso (aumentando la densità dell'acqua) immette reattività positiva.

\subsubsection{\textbf{Meccanismi di feedback esogeni (sistemi di Controllo})}

\textbf{Barre di controllo} L'inserzione nel core di materiali assorbitori aumenta localmente le perdite per cattura. Nei modelli a parametri concentrati, questo effetto geometrico viene descritto tramite un termine di reattività globale dipendente dall'altezza di inserimento delle barre.

\subsection{Neutroni ritardati e precursori}

 Sebbene la maggioranza dei neutroni venga emessa quasi istantaneamente all'atto della fissione, una piccola frazione $\beta$ (circa lo 0,65\%) viene rilasciata in ritardo dal decadimento radioattivo di specifici frammenti di fissione, definiti precursori. Poiché questi isotopi decadono con tempi caratteristici che spaziano da frazioni di secondo a quasi un minuto, la fenomenologia viene modellata raggruppandoli in sei famiglie temporali principali. Sebbene la frazione $\beta$ sia numericamente esigua, i neutroni ritardati sono il perno della controllabilità dell'impianto: essi dilatano l'inerzia dinamica del reattore, portando i tempi di escursione della potenza da frazioni di millisecondo a decine di secondi, rendendo così materialmente possibile l'intervento dei sistemi meccanici di regolazione.

\section{Modello matematico}

Il comportamento neutronico del reattore è stato modellato attraverso le equazioni di cinetica puntuale con neutroni ritardati, assumendo che la distribuzione spaziale del flusso neutronico rimanga invariata e che l’evoluzione temporale possa essere descritta attraverso la sola variazione della densità neutronica. In tale approccio la dinamica del sistema è governata da un’equazione differenziale per la densità neutronica $n(t)$ e da un insieme di equazioni per le concentrazioni $C_i(t)$ dei precursori dei neutroni ritardati.

\begin{equation}
\frac{dn(t)}{dt} = \frac{\rho - \beta}{\Lambda} \, n(t) + \sum_{i=1}^{6} \lambda_i C_i(t)
\label {3.1}
\end{equation}
L’evoluzione delle concentrazioni dei precursori appartenenti ai sei gruppi di neutroni ritardati è descritta invece da:
\begin{equation}
\frac{dC_i(t)}{dt} = \frac{\beta_i}{\Lambda} \, n(t) - \lambda_i C_i(t)
\qquad i = 1, \dots, 6
\label {3.2}
\end{equation}
dove $\rho$  rappresenta la reattività del sistema, $\beta$ la frazione totale di neutroni ritardati, $\beta_i$ la frazione associata all’i-esimo gruppo, $\lambda_i$ la costante di decadimento dei relativi precursori e $\Lambda$ il tempo di vita medio dei neutroni pronti. Tali parametri sono introdotti nel codice come costanti numeriche basate su valori tipici per reattori termici.\\ \\La reattività è stata modellata considerando tre contributi di retroazione: temperatura del combustibile, temperatura del moderatore e posizione delle barre di controllo. Essa è espressa come combinazione lineare delle deviazioni rispetto alle condizioni di riferimento, usando la seguente equazione:
\begin{equation}
\rho = \alpha_f (T_f - T_{f0}) + \alpha_m (T_c - T_{c0}) + \alpha_h (h - h_0)
\end{equation}
dove $\alpha_f$, $\alpha_m$, $\alpha_h$ sono i coefficienti di reattività rispettivamente legati alla temperatura del combustibile, alla temperatura del moderatore e all'altezza delle barre di controllo. Questi coefficienti vengono assunti costanti nell'intervallo termico esplorato, introducendo un'approssimazione del primo ordine. Tale linearizzazione è accettabile per transitori operativi che non si discostano drasticamente dalle condizioni nominali, permette di descrivere accuratamente la risposta dinamica del sistema mantenendo un'elevata efficienza di calcolo. 
Nel modello dinamico implementato, la reattività totale $\rho$ è espressa mediante la scomposizione tra i contributi intrinseci di retroazione termica del nocciolo e l'azione esogena dei sistemi di controllo. Più specificamente, le variazioni termiche del combustibile e del moderatore costituiscono i feedback intrinseci stabili del sistema, mentre lo spostamento delle barre rappresenta l'input di controllo esogeno. \\In questa trattazione sono considerate delle temperature di riferimento pari alle temperature di progetto del reattore, ossia di temperatura media del combustibile pari a 550 °C e del moderatore di 310 °C. Questi valori sono poi stati modificati nel paragrafo \ref{sec: steadystate} in cui i valori dello stato stazionario sono calcolati attraverso le equazioni della termoidraulica del reattore. \\Le concentrazioni iniziali delle famiglie di precursori, sono calcolate, imponendo lo stato stazionario nell'equazione (\ref{3.2})

\begin{equation}
C_{i}(0) = \frac{\beta_i}{\lambda_i \Lambda} n(0)
\end{equation}

\textbf{In seguito, una piccola legenda sulle convenzioni dei dati utilizzati:}

\begin{itemize}
\item $n(t)$: densità neutronica
\item $C_i(t)$: concentrazione dei precursori del gruppo $i$
\item $\rho$: reattività
\item $\beta$: frazione totale di neutroni ritardati
\item $\beta_i$: frazione del gruppo $i$ di neutroni ritardati
\item $\lambda_i$: costante di decadimento dei precursori
\item $\Lambda$: Tempo medio di vita dei neutroni pronti
\item $T_f$, $T_c$: temperatura del combustibile e del moderatore
\item $h$: posizione delle barre di controllo
\end{itemize}

\section{Risultati}

Al fine di validare il modello e osservare la dinamica del nocciolo, il sistema di equazioni differenziali è stato integrato numericamente partendo dalle condizioni di equilibrio stazionario ($T_{f0}$ = 550\ °C, $T_{c0}$ = 310\ °C, $n_0$= 1.0). Per testare la stabilità del sistema, sono state imposte delle perturbazioni a gradino sui parametri di feedback termico ($T_f$, $T_c$) e meccanico ($h$), analizzando la conseguente evoluzione temporale di $n(t)$. 
\begin{figure}[ht!]
    \centering
    \begin{minipage}{0.48\textwidth}
        \centering
        \includegraphics[width=\linewidth]{im1.jpeg}
        \caption{$T_f=600 {^\circ} C, T_c=305 {^\circ} C, h=0 cm$}
        \label{fig:img1a}
    \end{minipage}
    \hfill
    \begin{minipage}{0.48\textwidth}
        \centering
        \includegraphics[width=\linewidth]{Plot_rods.jpeg}
        \caption{$T_f=550 {^\circ} C, T_c=310 {^\circ} C, h=1 cm$}
        \label{fig:img2a}
    \end{minipage}

    \vspace{0.5cm} 
    \begin{minipage}{0.6\textwidth}
        \centering
        \includegraphics[width=\linewidth]{im3.jpeg}
        \caption{$T_f=660 {^\circ}C, T_c=240 {^\circ}C, h=0$}
        \label{fig:img3a}
    \end{minipage}
\end{figure}\\
Nella prima figura \ref{fig:img1a} un aumento di 50 °C della temperatura del combustibile e una diminuzione di 5 °C della temperatura del moderatore portano ad una diminuzione della densità neutronica: l'effetto Doppler ha un impatto maggiore rispetto all'aumento di densità del moderatore.\\ Nella seconda figura \ref{fig:img2a}, a un abbassamento delle barre di controllo consegue un naturale calo nella densità neutronica, dato dall'intervento esogeno sulla reazione a catena. \\Nell'ultimo grafico \ref{fig:img3a}, un aumento di 110°C nella temperatura del combustibile e una diminuzione di 70°C nella temperatura del moderatore, conducono a una crescita esponenziale della densità neutronica: questo avviene poiché il contributo neutronico dato dalla densificazione del moderatore supera di gran lunga lo smorzamento di reazione provocato dall'effetto Doppler.
 \\\\Avendo ora definito e validato la risposta puramente neutronica del reattore alle perturbazioni esterne, il prossimo passo consisterà nell'integrare queste equazioni all'interno della più ampia cornice termoidraulica. Nei capitoli successivi si procederà pertanto all'accoppiamento dei due modelli, al fine di catturare la reale natura interdipendente dei transitori operativi del reattore.

\chapter{Termoidraulica}
L’obiettivo del modello sviluppato è la costruzione di una rappresentazione termoidraulica semplificata, ma fisicamente consistente, del reattore. L’analisi è finalizzata allo studio del comportamento transitorio del sistema in seguito a una variazione di potenza, introdotta mediante un gradino positivo del 10\%, e alla valutazione della risposta dinamica del combustibile e del fluido refrigerante. Si è adottato un approccio a parametri concentrati, in cui il comportamento del sistema è descritto mediante i valori medi delle variabili di interesse, considerati uniformi nello spazio e dipendenti dal tempo. 

\section{Definizione del modello}
Per tener conto della dipendenza del coefficiente di scambio conduttivo e del calore specifico del combustibile, il nocciolo è stato diviso in 3 volumi concentrici, dove il volume interno è un cilindro pieno e i due più esterni sono cilindri anulari, ognuno con una propria temperatura e conduttività termica (dipendente dalla temperatura). Per la suddivisione dei volumi, i raggi scelti sono stati calcolati a partire dal raggio del fuel rod $r_f$. Nella zona centrale si ha $r < r_1$, nella zona intermedia si ha $r_1 < r < r_2$ e nella zona esterna si ha $r_2 < r < r_f$, con:
\begin{equation}
    r_1 = \frac{r_f}{\sqrt{3}}, \qquad r_2 = r_f\sqrt{\frac{2}{3}}
\end{equation}
\begin{figure}[H]
    \centering
    \includegraphics[width=0.4\linewidth]{Fuel rod.jpeg}
    \caption{Suddivisione dell'elemento combustibile}
    \label{fig:placeholder1}
\end{figure}
Il motivo di questa suddivisione risiede nella grande variabilità della temperatura all'interno del combustibile. Nonostante il diametro sia dell'ordine di qualche millimetro, la variazione tra il centro e la superficie del combustibile è di circa 300 °C.\\ \\La resistenza termica totale è stata modellata come somma in serie dei contributi di conduzione e convezione:
\begin{equation}
    R_{tot} = R_{fuel} + R_{gap} + R_{clad} + R_{conv}
\end{equation}
con:
\begin{equation}    
    R_{fuel} = \frac{1}{4 \pi H k_f(T_f)},
\end{equation}
\begin{equation}
    R_{gap} = \frac{1}{2 \pi r_f H h_{gap}},
\end{equation}
\begin{equation}
    R_{clad} = \frac{\ln\left(\frac{r_{ce}}{r_{ci}}\right)}{2 \pi H k_{clad}},
\end{equation}
\begin{equation}
    R_{conv} = \frac{1}{2 \pi r_{ce} H h_c}.
\end{equation}
Le resistenze del gap di elio e del \textit{cladding} sono fissate: la prima con un coefficiente di scambio termico convettivo posto pari a $1593.77 \frac {W}{m^2\cdot K}$ mentre la seconda con una conduttività pari a $17 \frac{W}{m\cdot K}$. Infine, la resistenza convettiva dell'acqua dipende dalla geometria del reattore e dalle proprietà dell'acqua. In quest'ultimo caso è stato anche considerato il contributo del termine
\begin{equation}
    2 \cdot \dot m_\text{channel} \cdot c_{p,c} (T_c-T_\text{inlet})
\end{equation}
Il coefficiente di scambio convettivo $h_c$ è stato valutato tramite correlazione di Dittus-Boelter corretta per bundle di barre, in funzione dei numeri adimensionali di Reynolds e Prandtl:
\begin{equation}
    Nu = 0.023 \, Re^{0.8} \, Pr^{0.4} \, C_{weisman}
\end{equation}
\begin{equation}
    h_c = \frac{Nu \, k}{D_h}
\end{equation}
\begin{equation}
    Re = \frac{\rho v D_h}{\mu}
\end{equation}
\begin{equation}
    Pr = \frac{\mu c_p}{k}
\end{equation}
dove il diametro idraulico $D_h$
\begin{equation}
    D_h = D \left( \frac{2\sqrt{3}}{\pi} \left(\frac{P}{D}\right)^2 - 1 \right)
\end{equation}
è stato ricavato a partire dalla geometria del reticolo triangolare.\\
La capacità termica del combustibile è stata modellata come funzione della temperatura tramite la correlazione semi-empirica di Fink-Touloukian:
\begin{equation}
    c_f(T) = 4.186 \cdot  10^{-2}  (710 + 0.6 \cdot T+ 4 \cdot (\frac{T}{1000})^7)
\end{equation}
Anche il coefficiente di conduttività termica $k$ è espresso in funzione dalla temperatura:
\begin{equation}
    k = \frac{100}{11.8+0.0238 \cdot T}+ 8.775 \cdot 10^{-15}T^3
\end{equation}
Il modello si concentra sulla dinamica termica del combustibile, descritta attraverso un bilancio energetico non stazionario. In particolare, la variazione temporale della temperatura del combustibile è data dalla differenza tra la potenza generata per fissione e il calore trasferito al fluido refrigerante.

\section{Condizioni iniziali}
\label {sec: steadystate}
Prima di osservare l'andamento delle grandezze d'interesse abbiamo cercato le effettive temperature in stato stazionario per il nostro modello, per avere una maggiore accuratezza sui valori iniziali del transitorio.\\ \\I valori di temperatura in condizioni stazionarie sono stati determinati mediante bilanci energetici. Nota la potenza termica totale generata dal \textit{fuel}, pari a $175 \text{MW}$, e ripartita tra le tre zone tramite una media pesata sui rispettivi volumi \eqref{eq:qi}, tale potenza deve eguagliare quella scambiata, espressa come rapporto tra la differenza di temperatura e la corrispondente resistenza termica.\\\\
In regime stazionario le derivate temporali si annullano, e il sistema di ODE si riduce a un sistema algebrico non lineare: per ciascun nodo la potenza generata o entrante (o entrambe) deve eguagliare quella uscente. La potenza generata nel volume $i$-esimo del combustibile è data da:
\begin{equation}
    \dot{Q}_i = P_\text{rod}\,\frac{V_i}{V_1+V_2+V_3}, \qquad i=1,2,3
    \label{eq:qi}
\end{equation}
Per il volume centrale si ha:
\begin{equation}
    \dot{Q}_1 = \dot{Q}_{12} \quad\Longrightarrow\quad
    P_\text{rod}\,\frac{V_1}{V_1+V_2+V_3} = \frac{T_1 - T_2}{R_{12}(T_1)}
    \label{eq:ss_nodo1}
\end{equation}
In maniera analoga si trovano i bilanci per il volume intermedio, per quello esterno, per il cladding e per il coolant, dove il termine di uscita è la portata calcolata come: 
$2\,\dot{m}_{ch}\,c_{p,c}(T_c - T_{in})$. Il sistema risultante viene risolto numericamente, fornendo come punto di partenza una stima iniziale accurata, ottenuta tramite il modello a resistenza totale concentrata.

\section{Transitorio}
Per il combustibile, il bilancio energetico è espresso come:
\begin{equation}
    m_fc_{pf}(T_f)\,\frac{dT_f}{dt} = P(t) - \frac{T_f - T_c}{R_{tot}(T_f)}
\end{equation}
dove $m_f$ è la massa di combustibile, $c_{pf}(T_f)$ il calore specifico dipendente dalla temperatura, $P(t)$ la potenza termica generata per barra e $R_{tot}$ la resistenza termica equivalente tra combustibile e refrigerante.\\
Per il refrigerante invece, il bilancio energetico assume la forma:
\begin{equation}
    m_c c_{pc} \,\frac{dT_c}{dt} = \frac{T_f - T_c}{R_{tot}} - 2 \dot{m} c_p (T_c - T_{in})
\end{equation}
dove $m_c$ è la massa di fluido nel canale, $c_{pc}$ il calore specifico, $\dot{m}$ la portata massica per canale e $T_{in}$ la temperatura di ingresso del refrigerante.\\ \\Il transitorio simulato consiste in un incremento a gradino della potenza del 10\% in corrispondenza del tempo $t = 10\ \si{\second}$:

\begin{equation}
    P(t) = \begin{cases} 
        P_0 & t \leq 10\ \si{\second} \\ 
        1.1 \cdot P_0 & t > 10\ \si{\second} 
    \end{cases}
\end{equation}
L’approccio adottato consente di catturare la dinamica termica accoppiata combustibile–refrigerante, risultando particolarmente adatto per analisi transitorie rapide e studi parametrici.\\ \\Le immagini seguenti riportano i risultati, permettendo di osservare l’evoluzione della temperatura del combustibile e del refrigerante nel tempo, e di valutare la rapidità e l’intensità della risposta termica del sistema.

\begin{figure}[!ht]
    \centering
    \includegraphics[width=1\linewidth]{TH1.png}
    \caption{Andamento della temperatura al centro del combustibile ($T_f$) in seguito ad una aumento del 10\% della potenza termica}
    \label{fig:placeholder2}
\end{figure}

\begin{figure}[!ht]
    \centering
    \includegraphics[width=1\linewidth]{TH2.png}
    \caption{Andamento della temperatura sulla superficie del combustibile in seguito ad una aumento del 10\% della potenza termica}
    \label{fig:placeholder3}
\end{figure}

\begin{figure}[!ht]
    \centering
    \includegraphics[width=1\linewidth]{TH3.png}
    \caption{Andamento della temperatura del coolant in seguito ad una aumento del 10\% della potenza termica}
    \label{fig:placeholder4}
\end{figure}

\clearpage
\section{Neutronica e termoidraulica}

La fase finale del lavoro riguarda l'accoppiamento tra la neutronica e la termoidraulica, volto all'analisi della propagazione degli effetti dovuti a un transitorio originato in uno dei due sottosistemi sulle grandezze caratteristiche dell'altro.\\ \\Nelle sezioni precedenti i due aspetti sono stati trattati separatamente: la risposta neutronica a una variazione della posizione delle barre di controllo, e la risposta termica del combustibile a un incremento a gradino della potenza generata. In questa sezione i due transitori sono invece simulati in modo accoppiato: la posizione delle barre di controllo, assegnata dall'utente come condizione iniziale a $t = 0$ s, introduce una perturbazione di reattività che modifica la potenza neutronica; a $t = 30$ s si sovrappone un aumento a gradino del $10\%$ della portata massica del primario:
\begin{equation}
    \dot{m}(t) =
    \begin{cases}
        \dot{m}_0        & t < 30\ \si{\second} \\
        1.1\,\dot{m}_0  & t \geq 30\ \si{\second}
    \end{cases}
\end{equation}
Questo secondo transitorio raffredda il coolant primario, riducendo $T_p$ e introducendo reattività positiva attraverso il coefficiente $\alpha_m$, che tende ad aumentare la potenza neutronica. L'accoppiamento è realizzato introducendo nella reattività la dipendenza dalle temperature calcolate dal modello termico:
\begin{equation}
    \rho(t) = \alpha_f\bigl(T_2(t) - T_{2,0}\bigr)
            + \alpha_m\bigl(T_p(t) - T_{p,0}\bigr)
            + \alpha_h\bigl(h - h_0\bigr)
    \label{eq:reattivita}
\end{equation}
dove $T_2$ è la temperatura del nodo intermedio del combustibile, $T_p$ è la temperatura media del coolant e $h$ è la posizione delle barre di controllo assegnata a $t = 0$ s. Sostituendo \eqref{eq:reattivita} nelle equazioni della cinetica puntiforme \eqref{3.1}--\eqref{3.2}, e accoppiando queste alle equazioni termiche del combustibile tramite il termine sorgente:
\begin{equation}
    q_i(t) = \frac{P_0\,n(t)\,V_i}{V_1+V_2+V_3}, \qquad i=1,2,3
\end{equation}
si ottiene un sistema di 12 equazioni differenziali ordinarie accoppiate: sei per i precursori, una per la densità neutronica e cinque per le temperature. La risposta integrata del sistema viene riportata nei grafici seguenti:

\begin{figure}[H]
    \centering
    \begin{minipage}{1\textwidth}
        \centering
        \includegraphics[width=0.9\linewidth]{image1nth.png}
        \label{fig:img1}
    \end{minipage}
    \hfill
    \begin{minipage}{1\textwidth}
        \centering
        \includegraphics[width=0.9\linewidth]{image2nth.png}
        \caption{Aumento di portata: $10\%$, inserimento control rods: $h=0 cm$}
        \label{fig:img2}
    \end{minipage}
\end{figure}

\begin{figure}[H]
    \centering
    \begin{minipage}{1\textwidth}
        \centering
        \includegraphics[width=0.9\linewidth]{image3nth.png}
        \label{fig:img1}
    \end{minipage}
    \hfill
    \begin{minipage}{1\textwidth}
        \centering
        \includegraphics[width=0.9\linewidth]{image4nth.png}
        \caption{Aumento di portata: $10\%$, inserimento control rods: $h=2 cm$}
        \label{fig:img2}
    \end{minipage}
\end{figure}


\chapter{Modello numerico dello Steam Generator a cassettoni}
\label{chap:scambiatore}

\section{Introduzione}
\label{sec:intro_scambiatore}

Dopo aver modellato la termoidraulica del core, in questa sezione si analizzerà il comportamento dello scambiatore di calore. Esso permette di trasferire la potenza termica generata grazie alle reazioni di fissione dal fluido primario al fluido secondario. Poiché il nostro reattore, essendo un SMR, non produce una grande quantità di calore, è necessario che lo scambiatore sia molto efficiente per ottimizzare la potenza elettrica utilizzabile. Il modello calcola i profili di temperatura ed entalpia del fluido primario e secondario lungo la direzione assiale dello scambiatore, sia in condizioni di regime stazionario sia dopo un transitorio operativo indotto da una variazione della portata secondaria definita dall'utente.\\
L'obiettivo della simulazione, in particolare, è duplice:
\begin{itemize}
    \item determinare lo stato stazionario del sistema (campo di temperatura del primario e entalpia e temperatura del secondario) a partire dalle condizioni di ingresso di pressione e temperatura/entalpia dei due fluidi;
    \item valutare la risposta dinamica del sistema a un cambiamento della portata massica secondaria regolata dalla turbina, quantificando la nuova condizione di equilibrio e mostrando il comportamento in transitorio sia della temperatura di uscita del primario sia della potenza scambiata nel tempo.
\end{itemize}
Nella nostra trattazione sono presenti 4 scambiatori di calore di tipologia once-through a tubi concentrici. La modellazione delle resistenze dello scambiatore di calore è analoga a quella adottata per le resistenze del core; è stata tuttavia trascurata la parte iniziale dello scambiatore poiché la quantità di potenza termica scambiata in quel tratto è trascurabile rispetto a quella scambiata nel resto del sistema, come si vedrà successivamente.

\begin{center}
\begin{table}[H]
    \centering
    \caption{Valori geometrici scambiatore di calore}
    \label{4.1}

\begin {tabular}{|m{8cm}|c|c|}
    \hline
    
    Grandezza & Cassetti & Sezioni\\
    \hline
    \hline
    Sezione trasversale tubi di scambio & $13 \times 1.5 mm$ & $8 \times 1.5 mm$\\
    \hline
    Numero cassetti per sezioni & 3 & -\\
    \hline
    Numero di moduli & 7 & 21\\
    \hline
    Numero di tubi di scambio termico nei moduli per sezioni & \multicolumn{2}{c|}{118} \\
    \hline
    Numero di tubi di scambio termico & 826 & 2478\\
    \hline
    Materiale dei tubi & \multicolumn{2}{c|}{Lega di PT-7M}   \\
    \hline
    Superficie di scambio termico & 93.4 $m^2$ & 280 $m^2$\\
    \hline
    Potenza termica & 14.5 MW & 43.7 MW\\
    \hline

\end{tabular}
\end{table}
\end {center}



\section{Geometria del problema e ipotesi di simmetria}
\label{sec:geometria}

Lo scambiatore è costituito da tre canali concentrici, ciascuno con simmetria cilindrica rispetto all'asse longitudinale $z$ del tubo (Figura~\ref{6.2}):
\begin{itemize}
    \item il \textbf{tubo interno}, in cui scorre una parte del fluido primario;
    \item l'\textbf{annulus}, la corona circolare compresa tra il tubo interno e il tubo esterno, in cui scorre il fluido secondario;
    \item il \textbf{mantello} (\emph{shell}), esterno al tubo esterno, in cui scorre la restante parte del fluido primario, con i tubi disposti secondo un passo triangolare.
\end{itemize}
\begin{figure}[H]
    \centering
    \includegraphics[width=0.5\linewidth]{SG2.png}
    \caption{}
    \label{6.2}
\end{figure}
\begin{figure}[H]
    \centering
    \includegraphics[width=1\linewidth]{SG1.png}
    \caption{Suddivisione dell'elemento combustibile}
    \label{6.1}
\end{figure}
La simmetria cilindrica delle tre regioni permette due semplificazioni fondamentali, su cui si fonda l'intero modello:
\begin{enumerate}
    \item \textbf{Riduzione dimensionale (3D $\rightarrow$ 1D).} Grazie alla simmetria cilindrica, tutte le grandezze di interesse (temperatura, entalpia, portata) possono essere mediate su ciascuna sezione trasversale e trattate come funzioni della sola coordinata assiale $z\in[0,L]$, dove $L$ è la lunghezza attiva dello scambiatore. Il problema, originariamente tridimensionale, si riduce così a un sistema di equazioni differenziali ordinarie (ODE) lungo $z$.
    \item \textbf{Resistenze di conduzione in geometria cilindrica.} La conduzione termica attraverso le pareti tubolari (di spessore finito) è descritta dalla classica espressione della resistenza conduttiva per un guscio cilindrico di lunghezza $dz$:
    \begin{equation}
        R_{\text{cond}} = \frac{\ln(r_{\text{ext}}/r_{\text{int}})}{2\pi \, k \, dz}
        \label{eq:R_cond_cilindrica}
    \end{equation}
    dove $r_{\text{int}}$ ed $r_{\text{ext}}$ sono i raggi interno ed esterno del tubo e $k$ è la conducibilità termica del materiale. Questa espressione, valida esclusivamente in geometria cilindrica viene applicata sia alla parete del tubo interno sia a quella del tubo esterno. La geometria è caratterizzata dai diametri $(D_{1,int},D_{1,ext}) = (5,8)\text{ mm}$ per il tubo interno e $(D_{2,int}, D_{2,ext}) = (10, 13)\text{ mm}$ per il tubo esterno, con un passo triangolare del reticolo pari a: $$P = 1.09 \cdot D_{2,ext}$$
    Ulteriori ipotesi sono legate inoltre alla costanza nella proprietà nei materiali e il trascuramento delle variazioni di volume del fluido.
\end{enumerate}
Il sistema viene discretizzato in $N$ celle di volume finito di lunghezza $dz = L/N$, ciascuna trattata come un volume di controllo con proprietà uniformi. I due fluidi scorrono in \textbf{controcorrente}: il primario entra a $z=L$ (caldo) ed esce a $z=0$ (freddo); il secondario entra a $z=0$ (subcooled) ed esce a $z=L$ (vapore surriscaldato). Questa configurazione, scelta per massimizzare lo scambio termico, impone una direzione di avanzamento opposta per i due fluidi nella risoluzione numerica.\\
Infine, l'impianto reale è costituito da un numero elevato di sotto-canali identici disposti in parallelo ($n_{\text{sub}}$, dato dal prodotto tra numero di gruppi, tubi per gruppo e sezioni); il modello simula un singolo sotto-canale, assumendo che la portata totale si distribuisca equamente tra tutti i canali, dove il generico sotto-canale è stato discretizzato in $N = 1000$ celle di uguale lunghezza $\Delta z = L/N$, con $L = 1.9\text{ m}$, secondo uno schema ai volumi finiti. Il numero totale di sotto-canali paralleli per l'intero sistema è $n_{sub} = 118 \times 7 \times 12 = 9912$. La portata totale del primario ($\text{flowrate I} = 902.8\text{ kg/s}$) e del secondario ($\text{flowrate II} = 77.5\text{ kg/s}$) vengono equamente ripartite tra i sotto-canali, definendo le portate locali $\dot{m}_I$ e $\dot{m}_{ann}$. La sezione del primario viene ulteriormente suddivisa tra l'area interna e l'area del mantello:  
\begin{equation}
    A_{\text{int}} = \pi \left( \frac{D_{1,\text{int}}}{2} \right)^2,
\end{equation}
\begin{equation}
    A_{\text{mant}} = \frac{\sqrt{3}}{2} \left( 1.09 \cdot D_{2,\text{ext}} \right)^2 - \frac{\pi D_{2,\text{ext}}^2}{4}
\end{equation}
proporzionalmente alle rispettive aree di passaggio del flusso. La geometria di tutto lo steam generator e delle varie sezioni può essere apprezzata nelle figure \ref {6.2} e \ref {6.1}.


\section{Resistenze termiche e correlazioni convettive}
\label{sec:resistenze}

Per ciascuna cella assiale, lo scambio termico tra primario e secondario avviene attraverso un percorso di resistenze in serie: convezione lato primario $\rightarrow$ conduzione nella parete del tubo $\rightarrow$ convezione lato secondario. Poiché il primario lambisce sia il tubo interno che il tubo esterno, esistono \emph{due} percorsi di scambio paralleli verso l'\textit{annulus}, schematizzati in Figura~\ref{sg3}.


\begin{figure}[H]
    \centering
    \includegraphics[width=0.7\linewidth]{immagine per andre.png}
    \caption{Schema a resistenze termiche in serie per una singola cella, lungo i due percorsi (tubo interno e tubo esterno)}
    \label{sg3}
\end{figure}

\subsection{Coefficienti di scambio convettivo}

Il coefficiente di scambio termico convettivo $h$ è stimato tramite la correlazione di \textbf{Dittus--Boelter}:
\begin{equation}
    Nu = 0.023\, Re^{0.8}\, Pr^{0.4}, \qquad h = \frac{Nu \cdot k_w}{D_h}
    \label{eq:dittus_boelter2}
\end{equation}
con $Re = \rho v D_h/\mu$ e $Pr=\mu c_p/k_w$, valida per flusso turbolento monofase completamente sviluppato.\\
Per il \textbf{tubo interno} e l'\textbf{\textit{annulus}}, dove la sezione di flusso è semplice (cerchio o corona circolare), si applica direttamente l'Eq.~\eqref{eq:dittus_boelter2} con il diametro idraulico geometrico ($D_h = D$ per il tubo, $D_h = D_{2,\text{int}}-D_{1,\text{ext}}$ per l'annulus).\\
Per il \textbf{mantello} invece, la presenza dei tubi disposti a passo triangolare modifica sia l'area di passaggio che il diametro idraulico:
\begin{equation}
    A_{\text{flow}} = \frac{\sqrt{3}}{2} P^2 - \frac{\pi D_{\text{ext}}^2}{4}, \qquad
    D_h = D_{\text{ext}}\left(\frac{2\sqrt{3}}{\pi}\left(\frac{P}{D_{\text{ext}}}\right)^2 - 1\right)
    \label{eq:area_mantello}
\end{equation}
dove $P$ è il passo triangolare tra i centri dei tubi adiacenti. Il numero di Nusselt viene inoltre corretto con un fattore geometrico empirico, detta correzione di Wiesman:
\begin{equation}
    Nu_{\text{mantello}} = 0.023\, Re^{0.8}\, Pr^{0.4}\left(1.106\,\frac{P}{D_{\text{ext}}} - 0.013\right)
    \label{eq:nu_mantello}
\end{equation}
che tiene conto dell'aumento di turbolenza indotto dalla disposizione triangolare dei tubi rispetto al caso di tubo isolato.


\subsection{Trattamento delle fasi del secondario}
\label{sub:fasi}

Il fluido secondario, durante l'attraversamento dello scambiatore, passa da liquido sottoraffreddato a vapore surriscaldato, transitando per una zona bifase. La fase locale viene determinata a partire dall'entalpia $h$ della cella confrontandola con i valori di saturazione $h_{L,\text{sat}}$ e $h_{V,\text{sat}}$ a pressione $P_{II}$ costante:
\begin{equation}
    \text{fase} =
    \begin{cases}
        \text{liquido} & h \le h_{L,\text{sat}} \\
        \text{bifase} & h_{L,\text{sat}} < h < h_{V,\text{sat}} \\
        \text{vapore} & h \ge h_{V,\text{sat}}
    \end{cases}
    \label{eq:fasi}
\end{equation}
Nella zona bifase, il titolo di vapore $x$ è calcolato con la regola della leva sull'entalpia:
\begin{equation}
    x = \frac{h - h_{L,\text{sat}}}{h_{V,\text{sat}} - h_{L,\text{sat}}}
    \label{eq:titolo}
\end{equation}
e la temperatura è fissata pari alla temperatura di saturazione $T_{\text{sat}}(P_{II})$, mentre nelle altre due regioni la temperatura è ricavata direttamente dalle tavole del vapore come $T = T(P_{II}, h)$.\\
Il coefficiente convettivo lato secondario dipende dalla fase:
\begin{itemize}
    \item in fase liquida o vapore si usano le correlazioni fornite da pyXsteam con le proprietà valutate, rispettivamente, tramite la temperatura e pressione del fluido.
    \item In zona bifase satura si adotta una formulazione semianalitica per l'ebollizione convettiva basata sul titolo termodinamico di vapore $x$, dove la correlazione risulta quella empirica data dalla correlazione di Saha. Per evitare discontinuità numeriche e modellare il fenomeno del dryout, il coefficiente di scambio termico $H$ viene differenziato in due regimi:
    \begin{equation}
        h_{\text{bif}} = h_l\left(1 + 3.8\left(\frac{x}{1-x}\right)^{0.76}\left(\frac{3.4}{p_{\text{crit}}}\right)^{0.38}\right), \qquad x<0.8
        \label{eq:saha}
    \end{equation}
    con $hl = 6000$, mentre per $x \ge 0.8$ (regime di post-dryout / mist flow): viene eseguita un'interpolazione lineare continua tra il picco calcolato a $x = 0.8$ ($H_{08}$) e il coefficiente convettivo del vapore surriscaldato ($H_{II,v}$):
     $$H = H_{08} + (H_{II,v} - H_{08}) \cdot \frac{x - 0.8}{1.0 - 0.8}$$
     \end{itemize}


\subsection{Resistenza totale di cella}

Combinando convezione primario, conduzione di parete (Eq.~\ref{eq:R_cond_cilindrica}) e convezione secondario, la resistenza totale lungo ciascuno dei due percorsi (tubo interno $i$, tubo esterno $e$) si sono create due funzioni dove i valori delle resistenze per la cella $k$-esima è:
\begin{equation}
    R_{i} = R_{c,i} + R_{\text{wall,int}} + \frac{1}{2\pi r_{1,\text{ext}}\, dz\, h_{II}}, \qquad
    R_{e} = R_{c,e} + R_{\text{wall,ext}} + \frac{1}{2\pi r_{2,\text{int}}\, dz\, h_{II}}
    \label{eq:R_totali}
\end{equation}
dove i valori di $R_{\text{wall}}$ generali, sempre ripresi da (Eq.~\ref{eq:R_cond_cilindrica}), sono dati da:
\begin{equation}
    R_{wall,i} = \frac{\ln(r_{i,ext}/r_{i,int})}{2\pi \cdot \Delta z \cdot k_{cond}}
\end{equation}
con conducibilità termica della parete pari a $k_{cond} = 12\text{ W m}^{-1}\text{ K}^{-1}$, mentre i generici $R_{c,i}$ e $R_{c,e}$ sono calcolati a ogni passo temporale in base ai rispettivi coefficienti locali di scambio termico e alle aree superficiali della cella con:
\begin {equation}
    \frac{1}{2\pi \cdot (D_{i,int}/2) \cdot \Delta z \cdot H_{int}}
\end{equation}
Si useranno, inoltre, nelle successive equazioni la conduttanza termica complessiva della cella, la quale è la somma delle due conduttanze in parallelo:

\begin{equation}
    G = \frac{1}{R_i} + \frac{1}{R_e}
    \label{eq:G_cella}
\end{equation}


\section{Equazioni discrete e schema di avanzamento temporale}
\label{sec:equazioni_discrete}

\subsection{Bilanci di energia sul volume di controllo}

Per ciascuna cella si scrive il bilancio di energia in forma di volume finito. Per il fluido primario nella cella $i$, l'energia accumulata varia per effetto del flusso advettivo entrante/uscente e del calore cedute al secondario:
\begin{equation}
    \rho_I A_{\text{tot},I}\, dz\, c_{p,I} \,\frac{T_i^{n+1}-T_i^{n}}{dt}
    = \dot m_I\, c_{p,I}\left(T_{i+1}^{n+1} - T_i^{n+1}\right) - G_i\left(T_i^{n+1} - T_{II,i}\right)
    \label{eq:bilancio_primario_continuo}
\end{equation}
dove il termine advettivo è scritto come un flusso \textit{upwind} rispetto alla direzione del flusso primario (da $i+1$ verso $i$, poiché il primario avanza da $z=L$ a $z=0$). Analogamente, per il secondario nella cella $i$:
\begin{equation}
    \rho_{s,i}\, A_{\text{ann}}\, dz\,\frac{h_i^{n+1}-h_i^{n}}{dt}
    = \dot m_{\text{ann}}\left(h_{i-1}^{n+1} - h_i^{n+1}\right) + \frac{Q_i}{1000}
    \label{eq:bilancio_secondario_continuo}
\end{equation}
con flusso \textit{upwind} da $i-1$ verso $i$ (il secondario avanza da $z=0$ a $z=L$) e $Q_i = G_i\,(T_i^{n+1}-T_{II,i})$ il calore scambiato nella cella (il fattore $1000$ converte da $\si{\watt}$ a $\si{\kilo\watt}$, dato che $h$ è espresso in $\si{\kilo\joule\per\kilo\gram}$).

\subsection{Discretizzazione temporale implicita}

L'avanzamento temporale è realizzato con uno schema di \textbf{Euler implicito} (\emph{backward}). La scelta è motivata dalla stabilità: uno schema esplicito sarebbe vincolato dalla condizione \textit{CFL}, $dt_{\max}\approx dz/v$, che per la discretizzazione assiale fine adottata impone passi temporali dell'ordine di $10^{-3}\,\si{\second}$, rendendo computazionalmente oneroso simulare transitori di decine di secondi. Lo schema implicito è invece incondizionatamente stabile, permettendo passi temporali ordini di grandezza più grandi senza insorgenza di oscillazioni numeriche.\\
Definendo i termini di accumulo (analoghi a una portata massica equivalente di cella):
\begin{equation}
    \alpha = \frac{\rho_I A_{\text{tot},I}\, dz}{dt}, \qquad
    \beta_i = \frac{\rho_{s,i}\, A_{\text{ann}}\, dz}{dt}
    \label{eq:alpha_beta}
\end{equation}
le Eq.~\eqref{eq:bilancio_primario_continuo}--\eqref{eq:bilancio_secondario_continuo}, risolte esplicitamente per l'incognita al passo $n+1$, divengono:
\begin{equation}
    T_i^{n+1} = \frac{\alpha\, T_i^{n} + \dot m_I\, T_{i+1}^{n+1} + \dfrac{G_i}{c_{p,I}}\, T_{II,i}}
    {\alpha + \dot m_I + \dfrac{G_i}{c_{p,I}}}
    \label{eq:T_implicita}
\end{equation}
\begin{equation}
    h_i^{n+1} = \frac{\beta_i\, h_i^{n} + \dot m_{\text{ann}}\, h_{i-1}^{n+1} + \dfrac{\delta Q_i}{1000}}
    {\beta_i + \dot m_{\text{ann}}}
    \label{eq:h_implicita}
\end{equation}
dove il termine di sorgente termica semimplicito è espresso da:
$$\delta Q_i = \max\left(0, T_I^{n+1}[i] - T_{II}^n[i]\right) \cdot G_i$$
Si tratta di un sistema di equazioni implicite poiché $T_{i+1}^{n+1}$ ed $h_{i-1}^{n+1}$ compaiono a loro volta come incognite. Tuttavia, grazie alla struttura unidirezionale dell'advezione upwind, il sistema non richiede l'inversione di una matrice: può essere risolto per \textbf{sostituzione sequenziale} (\emph{sweep}), seguendo l'ordine naturale imposto dal verso del flusso:
\begin{itemize}
    \item il \textbf{primario} viene aggiornato eseguendo lo sweep da $i=N-1$ a $i=0$: la cella $N-1$ usa come condizione al contorno $T_{I,\text{in}}$, e ogni cella successiva usa il valore $T^{n+1}_{i+1}$ appena calcolato;
    \item il \textbf{secondario} viene aggiornato con sweep da $i=0$ a $i=N-1$: la cella $0$ usa come condizione al contorno $h_{II,\text{in}}$, e ogni cella successiva usa $h^{n+1}_{i-1}$ appena calcolato.
\end{itemize}
Il termine di scambio $G_i(T_i-T_{II,i})$ è trattato in modo \textbf{semi-implicito}: la temperatura del secondario $T_{II,i}$ usata nell'aggiornamento del primario (Eq.~\ref{eq:T_implicita}) è quella nota dal passo precedente $n$, mentre il flusso $Q_i$ usato nell'aggiornamento del secondario (Eq.~\ref{eq:h_implicita}) impiega la temperatura del primario $T_i^{n+1}$ già aggiornata nello stesso passo. Questo accoppiamento sequenziale (primario $\to$ secondario) evita la soluzione di un sistema lineare $2N\times 2N$ completo, mantenendo comunque una buona stabilità grazie all'impiego implicito dei termini advettivi dominanti.\\\\
Un vincolo fisico aggiuntivo impone che il calore non possa fluire dal secondario al primario: se l'aggiornamento implicito restituirebbe $T_i^{n+1} < T_{II,i}$, la cella viene ricalcolata annullando il termine di scambio,
\begin{equation}
    T_i^{n+1} = \frac{\alpha\, T_i^{n} + \dot m_I\, T_{i+1}^{n+1}}{\alpha + \dot m_I}
    \label{eq:T_clamp}
\end{equation}
e analogamente il flusso $Q_i$ nell'Eq.~\eqref{eq:h_implicita} è limitato da $\max(0, T_i-T_{II,i})$.\\
Lo stato termodinamico, la temperatura $T_{II}$ e il titolo di vapore $x$ del secondario vengono aggiornati cella per cella interfacciando l'entalpia $h^{n+1}$ con le routine della libreria pyXSteam alla pressione operativa $P_{II} = 34\text{ bar}$, distinguendo i regimi di liquido sottoraffreddato ($h \le h_L^{sat}$), bifase saturo ($h_L^{sat} < h < h_V^{sat}$) e vapore surriscaldato ($h \ge h_V^{sat}$).

\section{Strategia di soluzione}
\label{sec:metodo_soluzione}

La soluzione completa è organizzata in tre fasi successive, ciascuna costruita sul risultato della precedente.

\subsection{Stato stazionario continuo: metodo di shooting}

La prima fase determina una stima della soluzione stazionaria tramite un \textbf{metodo di shooting}. Il problema stazionario, ottenuto annullando le derivate temporali nelle Eq.~\eqref{eq:bilancio_primario_continuo}--\eqref{eq:bilancio_secondario_continuo}, è un problema ai valori al contorno (la temperatura del primario è nota in $z=L$, ovvero $T_{p,i} = 313,5^\circ\text{C}$, e anche la temperatura del secondario in $z=0$, ovvero $T_{s,i} = 170^\circ\text{C}$, ma vanno propagati in direzioni opposte), dove le condizioni al contorno del sistema sono intrinsecamente accoppiate agli estremi opposti a causa della configurazione controcorrente. Si procede quindi con un algoritmo  che ipotizza per tentativi la temperatura di uscita del primario $T_{I,\text{out}}$ (in $z=0$): a partire da un valore di prova, si integra il sistema di ODE stazionarie
\begin{equation}
    \frac{dT_I}{dz} = \frac{G\,(T_I-T_{II})}{\dot m_I\, c_{p,I}}, \qquad
    \frac{dh_{II}}{dz} = \frac{G\,(T_I-T_{II})}{\dot m_{\text{ann}}}
    \label{eq:ode_stazionarie}
\end{equation}
con schema upwind esplicito da $z=0$ a $z=L$, e si verifica se la temperatura del primario integrata a $z=L$ coincide con la condizione al contorno nota $T_{I,\text{in}}$. La corretta condizione iniziale viene trovata minimizzando il residuo rispetto alla temperatura reale d'ingresso imposta ($T_{I,in} = 315.5^\circ\text{C}$) mediante il metodo di Brent (\textit{root scalar}).
Il profilo continuo così ottenuto viene proiettato sulla griglia spaziale tramite uno sfasamento geometrico (\textit{face-to-center correction}) per l'entalpia del secondario e utilizzato come input per un processo di warmup implicito a $\Delta t = 1\text{ s}$. Questo double-step numerico permette di convergere allo stato stazionario discreto in meno di 300 iterazioni, garantendo un perfetto bilancio energetico trilaterale (con uno squilibrio numerico tra i flussi del primario, del secondario e delle pareti inferiore allo $0.001\%$) prima dell'innesco del transitorio dinamico.


\subsection{Warmup implicito verso lo stazionario discreto}

La soluzione ottenuta dallo shooting è una soluzione accurata del problema \emph{continuo}, ma non coincide esattamente con la soluzione stazionaria del sistema \emph{discretizzato} a $N$ celle finite (gli schemi \textit{upwind} continuo e discreto introducono errori di troncamento leggermente diversi). Per ottenere una condizione iniziale consistente con lo schema numerico che verrà poi usato nel transitorio, il profilo dello \textit{shooting} viene fatto evolvere tramite lo schema implicito (Eq.~\ref{eq:T_implicita}--\ref{eq:h_implicita}) con un passo temporale grande ($dt=\SI{1}{\second}$), fino a convergenza. La convergenza è dichiarata quando la variazione massima di temperatura e di entalpia tra due passi successivi resta sotto tolleranza per un numero prestabilito di iterazioni consecutive, evitando falsi positivi dovuti a fluttuazioni numeriche transitorie.\\\\
Al termine del \textit{warmup}, la consistenza della soluzione viene verificata confrontando tre stime indipendenti della potenza termica scambiata complessivamente: la somma cella-per-cella dei flussi locali $Q=\sum_i G_i(T_i-T_{II,i})$, il bilancio energetico sul primario $Q = \dot m_I c_{p,I}(T_{I,\text{in}}-T_{I,\text{out}})$ e il bilancio sul secondario $Q = \dot m_{\text{ann}}(h_{\text{out}}-h_{\text{in}})$. La coincidenza (a meno di errori numerici minimi) tra le tre stime costituisce una verifica indipendente della correttezza dell'implementazione.

\subsection{Transitorio: salto di portata secondaria}

A partire dallo stato stazionario discreto, si simula un \textbf{aumento o diminuzione a crescita con saturazione} della portata massica secondaria all'istante $t=\SI{10}{\second}$, mantenendo invariate tutte le altre condizioni al contorno. Questo è stato fatto poiché si è ricercato di simulare il vero e proprio transitorio compiuto da una pompa nel sistema secondario. L'integrazione procede con lo schema implicito (Eq.~\ref{eq:T_implicita}--\ref{eq:h_implicita}) con passo temporale $dt=\SI{0.1}{\second}$ fino a $t=\SI{30}{\second}$, sufficiente a osservare sia il transitorio immediatamente successivo al salto sia il raggiungimento di un nuovo stato di equilibrio. A ogni passo, le resistenze convettive lato primario (Eq.~\ref{eq:dittus_boelter2}--\ref{eq:nu_mantello}) sono ricalcolate, poiché dipendono dalla portata tramite il numero di Reynolds, e quindi variano nel tempo in risposta al salto.\\
Le grandezze di interesse — temperatura di uscita del primario e potenza termica scambiata (calcolata con i tre metodi indipendenti) — vengono salvate a ogni passo temporale. Il confronto tra i valori medi nella finestra immediatamente precedente il salto e quelli nell'ultima finestra della simulazione fornisce la variazione netta indotta dal transitorio sulle grandezze di uscita.
\chapter{Andamento della Pressione}
Nel circuito secondario è presente una turbina che permette la regolazione della portata di vapore in uscita, che di conseguenza regola la potenza estratta. Questo influenza direttamente le condizioni termodinamiche all'interno dello scambiatore di calore. Un aumento della portata richiesta dalla turbina comporta un minor volume totale di vapore, quindi una diminuzione della pressione nel lato secondario. Tale riduzione di pressione favorisce un maggiore trasferimento di calore dal circuito primario al secondario, aumentando il tasso di generazione del vapore. Viceversa, una diminuzione della portata richiesta riduce il prelievo di vapore, causando un incremento della pressione nel lato secondario dello steam generator e una conseguente riduzione del gradiente termico disponibile per lo scambio di calore. Pertanto, la pressione nello steam generator è strettamente correlata alla portata del fluido secondario e rappresenta una delle principali variabili controllate per garantire il corretto bilanciamento energetico e la stabilità operativa dell'impianto.

\section{Modello approssimato}
Per il calcolo della variazione della pressione nel tempo, è stata adottata un'approssimazione del problema fisico reale tramite l'introduzione delle costanti $K_p$ e $tau_p$, che rappresentano rispettivamente la sensibilità della pressione alla portata e la costante di tempo. Il fluido secondario viene considerato sempre in condizioni di saturazione, il che è vero quando si trova in stato bifase ed è una buona approssimazione per le condizioni di uscita. La zona di liquido sottoraffreddato occupa una frazione limitata della lunghezza del generatore di vapore, poiché il salto entalpico necessario a portare il fluido alla saturazione è significativamente inferiore al calore latente di vaporizzazione. L'ipotesi $u \approx u_{sat}(P_s)$ introdotta nel bilancio di pressione introduce pertanto un errore limitato alla sola zona di ingresso del secondario, ed è ritenuta accettabile come approssimazione di primo ordine.\\ \\Il coefficiente $K_p$ del modello fenomenologico può essere derivato dal bilancio di energia interna sul volume di controllo del secondario. Indicando con $U = M_s\,u_{sat}(P_s)$ l'energia interna totale:
\begin{equation}
    \frac{dU}{dt} = \dot{Q}_{in} - \dot{m}_{II}\,\Delta h
\end{equation}
Poiché quest'ultima dipende dalla pressione, che dipende a sua volta dal tempo, si ha:
\begin{equation}
    \frac{dP_s}{dt} = \frac{\dot{Q}_{in} - \dot{m}_{II}\,\Delta h}
                           {M_s\,\dfrac{du_{sat}}{dP_s}}
    \label{eq:dPdt_fisica}
\end{equation}
In prossimità del regime stazionario, dove $\dot{Q}_{in,0} = \dot{m}_{II,0}\,\Delta h_0$, una perturbazione $\delta\dot{m}$ della portata secondaria produce:
\begin{equation}
    \frac{dP_s}{dt} \approx -\frac{\Delta h_0}{M_s\,\dfrac{du_{sat}}{dP_s}}
    \,\delta\dot{m} = -K_p\,\delta\dot{m}
\end{equation}
Il coefficiente $K_p$ risulta quindi:
\begin{equation}
    K_p = \frac{\Delta h_0}{M_s\,\dfrac{du_{sat}}{dP_s}}
    \qquad \left[\si{\bar\per(\kilogram\per\second\per\second)}\right]
\end{equation}
Il termine di richiamo $-(P_s - P_{s,0})/\tau_p$ nel modello fenomenologico rappresenta invece l'effetto stabilizzante del feedback termico: al diminuire di $P_s$, $T_{sat}$ scende, $\dot{Q}_{in}$ aumenta, e la pressione tende a stabilizzarsi al nuovo valore di equilibrio. La costante $\tau_p$ non emerge direttamente dal bilancio fisico ma parametrizza la velocità di questo feedback, che nel modello fisico \eqref{press2} è invece catturato implicitamente dalla dipendenza di $\dot{Q}_{in}$ da $P_s$ attraverso $T_{sat}(P_s)$.

\section{Risultati}

\begin {figure}[H]
    \centering
    \includegraphics[width=0.8\linewidth]{image1pr.png}
    \caption{Variazione nel tempo della pressione del secondario, dovuta ad un aumento della portata del secondario di +10\%}
    \label{press1}
\end{figure}

\begin {figure}[H]
    \centering
    \includegraphics[width=0.8\linewidth]{image2pr.png}
    \caption{Variazione nel tempo della pressione del secondario, dovuta ad un aumento della portata del secondario di -10\%}
    \label{press2}
\end{figure}

La costante di tempo della pressione è stata posta pari a quella della portata: $\tau_p = \tau_{fr} = 5\ \si{\second}$. È importante osservare che il tempo di risposta apparente della pressione nei grafici risulta maggiore di $\tau_p$: ciò è dovuto al fatto che il punto di equilibrio istantaneo non è fisso ma si sposta nel tempo insieme alla portata, che evolve con la propria costante di tempo. La pressione insegue quindi un obiettivo mobile, e la dinamica risultante produce una risposta apparentemente più lenta di $\tau_p$.\\ \\Il coefficiente $K_p$, calcolato a partire dal bilancio fisico illustrato, risulta essere $K_p = 52.36\times10^{-3}\ 
\si{\bar\per(\kilogram\per\second\per\second)}$.

\chapter{Analisi del sistema completo}

La simulazione del sistema completo rappresenta il risultato finale del lavoro, 
in cui i modelli sviluppati nei capitoli precedenti vengono integrati in un unico 
schema risolutivo accoppiato. L'obiettivo è riprodurre fedelmente la risposta 
dinamica del reattore RITM-200 a perturbazioni operative, tenendo conto dell'interdipendenza dei fenomeni neutronici, termoidraulici e termodinamici 
che governano il comportamento del sistema.

\section{Strategia di accoppiamento}

Il principale elemento di complessità nella simulazione del sistema completo 
risiede nell'accoppiamento tra il modello del nocciolo e quello dello scambiatore 
di calore. I due sottosistemi sono fisicamente connessi attraverso il circuito 
primario: la temperatura di uscita dal nocciolo (\textit{hot leg}) costituisce la 
condizione al contorno di ingresso per lo \textit{SG}, mentre la temperatura di 
uscita dallo \textit{SG} (\textit{cold leg}) determina la temperatura di ingresso 
del refrigerante nel nocciolo, influenzando direttamente la reattività attraverso 
il coefficiente di feedback del moderatore $\alpha_m$.\\\\
Questa dipendenza circolare — il nocciolo influenza lo \textit{SG} e lo 
\textit{SG} influenza a sua volta il nocciolo — rende il problema intrinsecamente 
accoppiato e richiede una strategia di soluzione che garantisca la consistenza 
fisica tra i due modelli ad ogni passo temporale.
La soluzione adottata consiste nell'integrare il sistema di ODE del nocciolo 
all'interno del loop temporale dello \textit{SG}, in modo che i due modelli 
condividano la stessa discretizzazione temporale. Ad ogni passo $\Delta t$:
\begin{enumerate}
    \item Lo \textit{SG} viene aggiornato con lo schema implicito a volumi finiti 
    (Eq.~\ref{eq:T_implicita}-\ref{eq:h_implicita}), usando come condizione al 
    contorno la temperatura di \textit{hot leg} calcolata al passo precedente.
    \item La temperatura di \textit{cold leg} aggiornata $T_{cold}^{n+1}$, 
    estratta dal profilo del primario all'uscita dello \textit{SG}, viene passata 
    al modello del nocciolo come condizione al contorno di ingresso.
    \item Il sistema di ODE del nocciolo — comprendente le equazioni della 
    cinetica puntiforme e il modello termoidraulico a tre nodi radiali — viene 
    integrato sull'intervallo $[t^n, t^{n+1}]$ tramite \texttt{solve\_ivp} con 
    metodo \texttt{LSODA}, producendo le temperature aggiornate del combustibile, 
    del rivestimento e del refrigerante, nonché la nuova densità neutronica $n^{n+1}$.
    \item La reattività viene ricalcolata con le temperature aggiornate secondo 
    l'Eq.~\eqref{eq:reattivita}, chiudendo il ciclo di feedback.
\end{enumerate}
Questo schema è classificabile come un metodo sequenziale: i due operatori fisici (trasporto nello \textit{SG} e cinetica del 
nocciolo) vengono risolti in sequenza all'interno dello stesso passo temporale, 
passandosi le condizioni al contorno aggiornate. L'approccio introduce un ritardo 
di un passo temporale nell'accoppiamento, il cui 
effetto sull'accuratezza della soluzione è trascurabile per i passi temporali 
adottati ($\Delta t = 0.2$~s), significativamente inferiori alle costanti di tempo 
dei transitori analizzati (dell'ordine di decine di secondi).

\section{Modello di pressione nel sistema accoppiato}

Nel sistema completo, la pressione del circuito secondario non è più un parametro 
fisso ma una variabile dinamica, aggiornata ad ogni passo temporale secondo il 
modello descritto nel Capitolo~5. La variazione di pressione modifica le proprietà 
di saturazione del fluido secondario (in particolare la temperatura di saturazione 
$T_{sat}(P_s)$ e le entalpie $h_L^{sat}$, $h_V^{sat}$) che vengono quindi 
ricalcolate ad ogni passo tramite le tavole \texttt{pyXSteam}. Questo introduce 
un ulteriore accoppiamento tra la dinamica della portata secondaria, la pressione 
e lo scambio termico nello SG, rendendo il sistema completamente 
retroazionato.

\section{Schema risolutivo}

Lo schema risolutivo complessivo della simulazione può essere sintetizzato nelle 
seguenti fasi:

\begin{enumerate}
    \item \textbf{Inizializzazione.} Dichiarazione dei parametri fisici e 
    geometrici del nocciolo e dello \textit{SG}; calcolo delle proprietà 
    termodinamiche nominali tramite \texttt{pyXSteam}.

    \item \textbf{Stato stazionario accoppiato.} Risoluzione iterativa del sistema 
    algebrico non lineare che accoppia le equazioni stazionarie del nocciolo 
    (risolte con \texttt{fsolve}) e il metodo di shooting per lo \textit{SG} 
    (con \texttt{root\_scalar}), fino alla convergenza della temperatura di 
    \textit{cold leg}.

    \item \textbf{Warmup accoppiato.} Integrazione implicita di 1000 passi con 
    $\Delta t_{wu} = 1$~s, volta a portare i profili discreti dello \textit{SG} 
    e del nocciolo allo stato stazionario consistente con lo schema numerico 
    adottato, eliminando le discontinuità numeriche dovute alla proiezione del 
    profilo continuo sulla griglia discreta.

    \item \textbf{Simulazione del transitorio.} Integrazione accoppiata 
    core-\textit{SG} con $\Delta t = 0.2$~s per la durata totale del transitorio, 
    con aggiornamento ad ogni passo della portata secondaria, della pressione, 
    delle resistenze convettive e della reattività.

    \item \textbf{Post-processing.} Salvataggio e visualizzazione delle variabili 
    di interesse: densità neutronica, temperature del combustibile e del 
    refrigerante, temperature di \textit{hot leg} e \textit{cold leg}, entalpia 
    e temperatura di uscita del secondario, potenza scambiata e pressione del 
    secondario.
\end{enumerate}



\chapter{Considerazioni finali sui risultati}
I risultati finali della simulazione mostrano i chiari effetti che gli scostamenti dalle condizioni nominali di portata del circuito secondario e della variazione di altezza delle barre di controllo producono sull'intero sistema, comprendente acqua refrigerante del circuito primario e fluido di lavoro del secondario.\\\\
Come primo caso di studio, viene esaminata una singola perturbazione consistente in un aumento del 10\% della portata nel circuito secondario.

\section{Variazione di portata}
\begin {figure}[H]
    \centering
    \includegraphics[width=1\linewidth]{image1fin.png}
    \caption{Steady state alle condizioni nominali, illustrante la distribuzione di temperatura del primario e del secondario lungo la lunghezza dello \textit{SGE}}
    \label{7.1}
\end{figure}
\begin {figure}[H]
    \centering
    \includegraphics[width=1\linewidth]{image2fin.png}
    \caption{Variazione nel tempo di: Portata del secondario, Pressione del secondario, Densità neutronica e Potenza termica scambiata}
    \label{7.2}
\end{figure}
Si osserva come gli andamenti delle variabili (ad esclusione della pressione per la quale è stata imposta una correlazione con la portata) presentino un andamento di natura esponenziale semplice simile. Questo accade per effetto di una relazione lineare tra portata massica e rispettivamente densità neutronica e potenza scambiata: 
\begin {equation}
        \dot Q = \dot m (T_{out}-T_{in}); \qquad \dot Q= \dot Q_0 \frac{n}{n_0}
\end {equation}
La pressione, invece, diminuisce conseguentemente all'aumento di vapore estratto dallo scambiatore, dato dall’incremento della portata nel circuito secondario che alimenta la turbina (e ne aumenta la potenza generata).
\begin {figure}[H]
    \centering
    \includegraphics[width=1\linewidth]{image3fin.png}
    \caption{Variazione nel tempo di: temperatura al centro del combustibile, temperatura sulla superficie del combustibile, coefficiente di scambio termico globale, temperatura del rivestimento}
    \label{7.3}
\end{figure}
Le temperature del combustibile crescono coerentemente con l’aumento della densità neutronica; 
Il coefficiente di scambio termico aumenta, invece, in accordo con l’incremento della portata del circuito secondario, e quindi del contributo del \textit{numero di Reynolds};
La temperatura del rivestimento decresce a causa della diminuzione di temperatura media del refrigerante (come mostrato in seguito).
\begin {figure}[H]
    \centering
    \includegraphics[width=1\linewidth]{image4fin.png}
    \caption{Variazione nel tempo di: Temperatura media del refrigerante, Temperatura di \textit{hot leg}, Temperatura di \textit{cold leg}, Temperatura in uscita del secondario}
    \label{7.4}
\end{figure}
L’andamento della temperatura del vapore nel circuito secondario è decrescente poiché la diminuzione della pressione comporta un abbassamento della temperatura di saturazione, dunque un surriscaldamento minore nel vapore. Il circuito primario segue esattamente l’andamento del circuito secondario, il quale asporta più calore e ne comporta un generoso raffreddamento.
\begin {figure}[H]
    \centering
    \includegraphics[width=1\linewidth]{image5fin.png}
    \caption{Variazione nel tempo di: Conduttanza media, Resistenza media lato Mantello, Resistenza media lato Tubo, Titolo medio del secondario}
    \label{7.5}
\end{figure}
L'andamento dei primi tre parametri è una diretta conseguenza dell’aumento del coefficiente di scambio termico, dovuto all’incremento della portata del secondario.
Il titolo medio, invece, diminuisce (sebbene non in modo significativo), siccome la maggior quantità di vapore estratta, per essere inviata in turbina, consente una prevalenza di fase liquida nello \textit{SGE} (il bilancio si assesta comunque su una massa di vapore non troppo inferiore rispetto a quella liquida). 
\begin {figure}[H]
    \centering
    \includegraphics[width=1\linewidth]{image6fin.png}
    \caption{Profili assiali di: Conduttanza totale, Resistenza a lato Mantello, Resistenza a lato Tubo, Titolo di vapore, Potenza locale scambiata}
    \label{7.6}
\end{figure}
La distribuzione lungo lo \textit{SGE} della Conduttanza mostra chiaramente la propensione allo scambio termico nettamente maggiore della fase liquida e bifase, rispetto al vapore che ha una conduttanza nettamente inferiore. Il fenomeno è osservabile anche attraverso gli andamenti delle resistenze; la situazione non varia significativamente allo stato finale, il quale presenta la sola differenza di avere una frazione maggiore di vapore bifase lungo l’asse.
Osservando il profilo della qualità del circuito secondario, nello stato finale, si nota il raggiungimento dell’unità appena in prossimità dell’estremo finale; ciò è indice del fatto che il vapore in ingresso in turbina subisce un surriscaldamento minimo, al contrario dello stato iniziale in cui il vapore ha più tempo per surriscaldarsi (essendo la portata minore).
Infine, in accordo con i coefficienti di scambio termico e con le differenze di temperatura tra primario e secondario, si apprezzano i profili (quasi invariati tra stato finale e iniziale) della Potenza scambiata lungo lo \textit{SGE}: la zona liquida concede un elevatissimo scambio termico all’estremo iniziale, dovuto per lo più alla grande differenza di temperatura, la quale tende a diminuire vertiginosamente con l’avvicinamento delle temperature; il profilo torna ad alzarsi con l’evoluzione dello stato bifase lungo l’asse; si conclude con un crollo dello scambio termico lungo la zona del vapore (come previsto, visto lo scarso coefficiente di scambio termico).
È interessante notare come lo stato bifase, grazie al fenomeno dell’ebollizione nucleata, mantenga alta la potenza scambiata lungo l’asse (al contrario dello stato liquido), nonostante il progressivo avvicinamento delle temperature. Questo è conseguenza della straordinaria capacità di scambio termico dei fluidi in ebollizione.

\section{Schema logico dei fenomeni analizzati nei grafici (1)}

L’obiettivo principale della simulazione è quello di capire in modo chiaro il comportamento dinamico del reattore a seguito di una possibile perturbazione in ambito operativo. 
Per ottenere ciò è fondamentale tenere conto del rapporto di causalità tra i fenomeni che si verificano all’interno del sistema.
L’aumento di potenza del reattore è giustificato da un incremento della richiesta di energia da parte della rete elettrica, che grava direttamente a carico della turbina a valle del secondario. È in questo punto che, tramite una valvola, si gestisce la portata di vapore in ingresso alla turbina; per un aumento di potenza la portata dovrà chiaramente aumentare (e.g. di +10\% ).\\ \\La maggiore portata provoca un aumento del calore asportato da parte del circuito secondario all’interno dello \textit{SG} e parallelamente una diminuzione di pressione (controllata dalla quantità di vapore presente nello \textit{SG}) nel lato secondario; questa diminuzione di pressione è controllata da una correlazione assunta lineare alla variazione della portata stessa.
L’aumento di portata e quindi di velocità di transito del secondario è responsabile dell’incremento generale del coefficiente di scambio termico, inoltre la diminuzione della pressione consente al secondario di posizionarsi ad una temperatura di saturazione inferiore, la quale aumenta la differenza di temperatura con il primario.
Gli incrementi di \textit{$\Delta$T} e \textit{U} permettono al circuito primario di fornire la maggiore potenza termica richiesta dal secondario.\\ \\ Nel primario, a cui viene sottratto più calore, consegue un calo di temperatura analogo a quello del circuito secondario, posizionandosi ad una temperatura di uscita dallo \textit{SG} (quindi di \textit{cold leg}) inferiore rispetto a quella nominale. 
Un abbassamento della temperatura di \textit{cold leg} si traduce in ingresso di acqua più fredda, e quindi più densa all’interno del core; data la duplice funzione di refrigerante e moderatore svolta dall’acqua del primario, un incremento di densità consente una migliore moderazione dell’attività neutronica e quindi un inserimento di reattività positivo all’interno del nocciolo.\\ \\Questo fenomeno di alta rilevanza è il motivo per il quale i reattori di tipo \textit{PWR} sono detti \textit{load followers}, ossia in grado di compensare automaticamente la maggiore richiesta di potenza da parte della turbina, attraverso un incremento naturale della reattività del nocciolo, il quale inevitabilmente sviluppa una potenza termica maggiore. Tutto ciò avviene quasi in contemporanea ai fenomeni termoidraulici. Questo si evince dal fatto che in 50 secondi tutte le variabili analizzate finiscono per convergere nel nuovo stato.\\ \\Infine, nonostante l’aumento della potenza termica nel core, la temperatura di \textit{hot leg} (in ingresso allo \textit{SG}) si stabilizza ad una temperatura inferiore alla precedente, provocando una discesa nella temperatura media del primario, mentre la temperatura del combustibile aumenta in accordo con la densità neutronica.\\ 
\section{Variazione delle barre di controllo}
È possibile analizzare il comportamento dovuto ad una variazione di altezza delle barre di controllo. Ad esempio, un inserimento di 3 cm di barre all’interno del \textit{core}.
I profili dello Steady state sono praticamente invariati
    

\begin {figure}[H]
    \centering
    \includegraphics[width=1\linewidth]{image7fin.png}
    \caption{Variazione nel tempo di: Portata del secondario, Pressione del secondario, Densità neutronica e Potenza termica scambiata
    }
    \label{7.7}
\end{figure}

La portata del secondario, e dunque la pressione, restano invariate. La densità neutronica, invece, diminuisce secondo un preciso andamento: il crollo iniziale è dovuto all’assorbimento immediato dei neutroni da parte delle barre di controllo, mentre la leggera risalita, che si stabilizza in un valore inferiore rispetto a quello iniziale, è dovuta al raffreddamento dell’acqua del circuito primario, con conseguente aumento della sua densità e della capacità moderante.\\ \\La potenza diminuisce, in quanto il calo di densità neutronica ostacola la produzione di potenza termica per fissione all’interno del combustibile.

\begin {figure}[H]
    \centering
    \includegraphics[width=1\linewidth]{image8fin.png}
    \caption{Variazione nel tempo di: Temperatura al Centro del Combustibile, Temperatura sulla Superficie del Combustibile, Coefficiente di scambio Globale, Temperatura del Rivestimento.}
    \label{7.8}
\end{figure}

La temperatura del combustibile (centro e superficie) segue un andamento molto simile a quello della densità neutronica in quanto è diretta conseguenza di quest’ultima. La temperatura del rivestimento, invece, diminuisce a seguito del raffreddamento (precedentemente introdotto) dell’acqua del circuito primario.\\ \\È interessante notare che, contrariamente alla diminuzione di potenza termica scambiata, il coefficiente di scambio termico tra il circuito primario e secondario aumenta. Ciò è dovuto alla presenza dominante di liquido e bifase all’interno dello \textit{SGE} (come verrà spiegato in seguito) che contribuiscono ad innalzare la qualità dello scambio termico.

\begin {figure}[H]
    \centering
    \includegraphics[width=1\linewidth]{image9fin.png}
    \caption{Variazione nel tempo di: Temperatura media refrigerante, Temperatura \textit{hot leg}, Temperatura \textit{cold leg}, Temperatura in uscita del secondario}
    \label{7.9}
\end{figure}

La diminuzione di temperatura dell’acqua del primario è conseguenza della diminuzione di temperatura del combustibile e della minore potenza termica assorbita. Il raffreddamento del circuito secondario, ugualmente, dipende dalla potenza scambiata con il primario, che trascina con sé il secondario nella discesa (ad esclusione del tratto della zona bifase, che permane a temperatura di saturazione)

\begin {figure}[H]
    \centering
    \includegraphics[width=1\linewidth]{image10fin.png}
    \caption{Variazione nel tempo di: Conduttanza media, Resistenza media lato Mantello, Resistenza media lato Tubo, Titolo medio del secondario.}
    \label{7.10}
\end{figure}

Gli andamenti della Conduttanza e della Resistenza (lato mantello e lato tubo) sono facilmente deducibili a seguito delle precedenti considerazioni riguardo il coefficiente di scambio termico.
La qualità media del secondario decresce. Questo è dovuto al crollo della quantità di vapore presente nello \textit{SGE} , causato dal raffreddamento generale, il quale lascia spazio alla presenza di liquido e bifase.

\begin {figure}[H]
    \centering
    \includegraphics[width=1\linewidth]{image11fin.png}
    \caption{Profili assiali di: Conduttanza totale, Resistenza a lato Mantello, Resistenza a lato Tubo, Titolo di vapore, Potenza locale scambiata}
    \label{7.1}
\end{figure}

L’analisi dei profili lungo l’asse risulta di particolare interesse: la conduttanza totale, come precedentemente discusso, aumenta nel tempo; infatti, si nota come nello stato finale rimane alta più a lungo rispetto allo stato iniziale, questo è dovuto unicamente alla mancanza di potenza termica entrante nell’acqua del circuito secondario, che fa più fatica ad evaporare e quindi sulla stessa lunghezza dello \textit{SGE} c’è una presenza maggiore di bifase. È chiaro che la conduttanza del bifase (nella zona centrale) e del liquido (nella zona sinistra) non ha motivo di variare tra i due stati, siccome la portata resta inalterata. Un risultato analogo si riscontra per le resistenze.
Così come nel caso precedente la qualità del secondario raggiunge l’unità in prossimità dell’estremo dello \textit{SGE} (sebbene per motivi differenti).
La potenza scambiata lungo l’asse ovviamente diminuisce globalmente nello stato finale, seppur mantenendosi alta più a lungo a causa della persistenza dello stato bifase.


\section{Schema logico dei fenomeni analizzati nei grafici (2)}

In questo scenario, la catena degli eventi risulta invertita rispetto al caso precedente. La causa principale della diminuzione di potenza si trova infatti a monte, nell’abbassamento delle barre di controllo, finalizzate appunto ad un controllo della potenza del nocciolo; risulta quindi un’azione diretta sulla variazione dello stato del reattore e non conseguenza della regolazione della turbina a valle come nel caso precedente.
L’inserimento delle barre provoca un immediato assorbimento di neutroni nel nocciolo e quindi un'inibizione nella reazione a catena di fissione, con conseguente calo di potenza generata.  Ciò provoca una diminuzione della temperatura del combustibile (diretta conseguenza), della temperatura del rivestimento e di quella del circuito primario (effetto della minore potenza transitante). È proprio la diminuzione di acqua del primario, e aumento della sua densità, che permette la leggera risalita nella densità neutronica.
Il paradosso interessante si riscontra nello \textit{SG}, nell’interazione termica tra primario e secondario. Fondamentalmente, il calore da scambiare diminuisce poiché così imposto dalle condizioni al contorno a monte; quindi, la temperatura di uscita del vapore secondario diminuisce. In questo modo, però, la presenza del bifase è dominante e il coefficiente di scambio termico aumenta. Il motivo per il quale il calore scambiato si assesta ad un valore inferiore è dovuto al fatto che la differenza di temperatura media tra primario e secondario diminuisce in modo molto più deciso rispetto all’aumento del coefficiente di scambio termico; ciò avviene perché il profilo del secondario è vincolato alla sua temperatura di saturazione, che mantiene la temperatura media pressoché invariata (al contrario di quella del primario, che diminuisce sensibilmente). Tale fenomeno non era riscontrato nel caso precedente, nel quale il vincolo della pressione veniva meno.
Chiaramente, il secondario, avendo la stessa portata, ma essendo più freddo, subirà un minor salto entalpico in turbina e la potenza erogata in rete sarà minore. 
In questo scenario la diminuzione della potenza erogata a valle è causata dalla catena di eventi che hanno inizio nel reattore.
Da questa analisi si distinguono due differenti modalità di parzializzazione del carico del reattore, adottate nei reattori di tipo \textit{PWR} (e, nel caso in esame, negli \textit{SMR}):

\begin{enumerate}
    \item \textbf{Controllo della potenza a monte.} Consiste in una riduzione della potenza, come analizzato in precedenza, poiché le barre di controllo risultano normalmente completamente estratte nelle condizioni di funzionamento nominali.

    Il comportamento del sistema segue in maniera coerente le regolazioni effettuate a monte, evolvendo di conseguenza e garantendo a valle il livello di potenza imposto dal sistema di controllo.

    \item \textbf{Controllo della potenza a valle.} Può comportare sia un aumento sia una diminuzione della potenza (in questa trattazione ne è stato analizzato il caso di un aumento) e sfrutta la capacità di \emph{load following} caratteristica dei reattori di tipo \textit{PWR}. La regolazione è effettuata agendo sulla turbina del circuito secondario per motivi di logistica e immediatezza. Il sistema si evolve, però, in modo tale da autoregolare la potenza che viene generata al nocciolo e renderla coerente con quella stabilita dalla turbina.
\end{enumerate}
La simulazione e analisi può anche avere come oggetto situazioni realistiche e più complesse, nelle quali le due azioni di controllo avvengono contemporaneamente ed è fondamentale per prevedere i comportamenti delle temperature dell’acqua, dei rivestimenti del combustibile e della reattività in modo tale da assicurarsi che rientrino nei limiti di sicurezza previsti per il corretto funzionamento dell’impianto.
\chapter*{Bibliografia}
\addcontentsline{toc}{chapter}{Bibliografia}
\printbibliography[heading=none]

\end{document}


